\section{Related work}
\label{sec:related_work}

\subsection{Financial dependencies and econophysics}

Econophysics studies financial markets using tools from statistical physics, nonlinear dynamics, complex systems and network science. A central empirical observation is that financial returns display stylized facts that are difficult to reconcile with independent Gaussian models. Heavy tails, volatility clustering, slow decay in absolute-return autocorrelations and collective market modes are recurring findings across markets and time scales \cite{mantegna_stanley_1999_book,cont_2001_empirical}. These patterns motivate empirical methods that focus on distributions, dependence and topology rather than only on mean returns.

Financial dependency analysis sits between econophysics and financial economics. Raddant and Di~Matteo \cite{raddant_dimatteo_2023} emphasize that correlations, causality measures, volatility spillovers and networks provide complementary views of market structure. The present paper follows this perspective. The empirical goal is not simply to estimate a correlation matrix, but to decompose it, filter it, transform it into networks and test whether the extracted structure has forecasting relevance.

Cross-correlation dynamics during crises have received particular attention. Zheng et~al. \cite{zheng_2012_changes} showed that changes in the cross-correlation structure can serve as early indicators of systemic instability. Billio et~al. \cite{billio_2012_econometric} proposed econometric measures of connectedness and systemic risk in financial sectors. Junior and Franca \cite{junior_2012_correlation} documented how correlations among global stock indices increase during crises, reducing the effective dimensionality of the market. These findings motivate the rolling-correlation and crisis-synchronization analysis undertaken in this paper.

\subsection{Random Matrix Theory in financial correlations}

RMT entered empirical finance as a way to separate information from noise in large correlation matrices. Laloux et~al. \cite{laloux_1999_noise,laloux_2000_random} showed that a large part of the eigenspectrum of financial correlation matrices can resemble the spectrum of random matrices described by the Mar\v{c}enko--Pastur law \cite{marcenko_1967_distribution}. Plerou et~al. \cite{plerou_2002_random} further studied eigenvalues and eigenvectors outside the random band and interpreted the leading eigenvalue as a market-wide mode. These studies provide the foundation for the spectral decomposition used here.

Subsequent work extended RMT-based cleaning methods for practical applications. Bun et~al. \cite{bun_2017_cleaning} reviewed and unified several approaches for cleaning large covariance matrices using random matrix theory, providing a comprehensive toolkit for practitioners and researchers. The mathematical foundations of random matrix theory are covered in the review by Guhr et~al. \cite{guhr_muller_1998}.

The RMT interpretation used in this manuscript is deliberately cautious. The largest eigenvalue is interpreted as the Market Mode because it captures broad positive co-movement across many assets. Eigenvalues above the Mar\v{c}enko--Pastur upper edge, excluding the first one, are interpreted as group or sector candidates. Eigenvalues inside the random band are treated as noisy residual components. This terminology is a filtering convention, not a claim that every eigenvector has a unique economic label.

\subsection{Financial networks: MST, PMFG and filtered graphs}

Mantegna \cite{mantegna_1999_hierarchical} introduced the MST as a way to reveal hierarchical organization in financial markets. The method converts correlations into distances and retains the connected tree with minimum total distance. Because the MST has only $N-1$ edges, it is extremely sparse. This sparsity is useful for visualization, but it also discards cycles and many local dependencies.

Subsequent studies applied the MST framework to study the dynamics and taxonomy of equity markets. Onnela et~al. \cite{onnela_2003_dynamics} analysed the dynamics of MST-based market correlations and their implications for portfolio analysis. Bonanno et~al. \cite{bonanno_2004_networks} extended the network approach to multiple world markets, showing that MSTs consistently organize assets by sector and geography. Song et~al. \cite{song_2012_evolution} studied the evolution of worldwide stock markets through correlation-based graphs, documenting how network structures change through time.

The PMFG relaxes the MST constraint while preserving planarity \cite{tumminello_2005_tool}. For $N$ nodes, it retains $3N-6$ edges and contains the MST as a subgraph when the same ranking of distances is used. The PMFG can preserve triangles and 4-cliques, making it more suitable for detecting local communities, sectoral clusters and hub reorganization. Tumminello et~al. \cite{tumminello_2010_correlation} further explored the relationship between correlation, hierarchies and networks in financial markets. More recently, Massara et~al. \cite{massara_2016_tmfg} introduced the Triangulated Maximally Filtered Graph (TMFG), which extends the PMFG framework and can be computed more efficiently for very large networks.

Filtered networks are particularly important for risk analysis. Pozzi et~al. \cite{pozzi_2013_spread} used PMFG-based analysis to show that peripheral assets in financial networks can offer better diversification than central hubs. Aste et~al. \cite{aste_2010_correlation} studied correlation dynamics in volatile markets using filtered network approaches. Kenett et~al. \cite{kenett_2012_dominating} used partial correlation analysis to reveal the dominating role of the financial sector in stock market networks. These studies motivate the hub-reorganization analysis presented in this paper, where we compare how centrality rankings change after RMT-based market-mode removal.

\subsection{Financial networks in emerging markets}

The application of network-based methods to emerging equity markets has attracted growing interest. Tabak et~al. \cite{tabak_2010_topology} studied the topology of the Brazilian interbank network, documenting its scale-free properties and vulnerability to targeted failures. Silva et~al. \cite{silva_2016_structure} analysed the structure and dynamics of the global financial network, with particular attention to how emerging economies connect to the broader system. Wang et~al. \cite{wang_2018_correlation} studied the correlation structure of world stock markets using Pearson and partial correlation-based networks, showing that emerging markets display distinctive dependency patterns.

Studies focused specifically on Brazilian markets include Pereira et~al. \cite{pereira_2019_multiscale}, who applied detrended cross-correlation analysis to construct multiscale networks of stock markets including B3, and Leonel Caetano and Yoneyama \cite{leonel_2022_brazilian}, who applied complex network analysis directly to the Brazilian stock market. Stosic et~al. \cite{stosic_2015_correlations} studied correlations and cross-correlations in Brazilian agrarian commodities and stocks. These studies demonstrate that Brazilian financial data present interesting structural properties, but a comprehensive pipeline combining RMT filtering, PMFG construction and volatility forecasting has not been systematically applied to a broad B3 equity universe.

\subsection{Volatility forecasting and machine learning}

Volatility forecasting has a long econometric tradition. GARCH-type models, introduced by Engle \cite{engle_1982_autoregressive} and generalized by Bollerslev \cite{bollerslev_1986_generalized}, capture conditional heteroskedasticity through parsimonious dynamic equations. Hansen and Lunde \cite{hansen_lunde_2005} showed that the simple GARCH(1,1) is notoriously difficult to outperform systematically. HAR-RV models \cite{corsi_2009_simple} approximate long-memory behaviour in realized volatility through daily, weekly and monthly components. The theoretical foundations of realized-volatility measurement were established by Andersen et~al. \cite{andersen_2003_modeling}. Poon and Granger \cite{poon_granger_2003} provided an influential review of volatility forecasting methods. Forecast evaluation is delicate because realized volatility is a proxy for latent conditional variance; QLIKE is often used because of its robustness properties under imperfect volatility proxies \cite{patton_2011_volatility}.

Machine-learning methods provide flexible nonlinear approximations. Random Forests \cite{breiman_2001_random}, gradient boosting \cite{friedman_2001_greedy} and regularized linear models \cite{hoerl_kennard_1970_ridge} can combine large feature sets containing lagged volatility, market-wide variables and network centralities. Bucci \cite{bucci_2020_realized} applied neural networks to realized volatility forecasting, showing that flexible models can capture nonlinear volatility dynamics. Christensen et~al. \cite{christensen_2022_realized} developed a machine learning framework for volatility forecasting and demonstrated the importance of careful feature engineering. Zhang and Broadstock \cite{zhang_2023_network_volatility} explored how financial network topology relates to market risk from a machine learning perspective.

In financial forecasting, however, flexibility can also amplify noise. For this reason, this paper evaluates machine learning against strong baselines and treats network features as potential incremental predictors rather than as a replacement for classical volatility dynamics. The nested feature-set design follows an ablation logic: Set~A measures the predictive value of classical volatility information, Set~B tests the incremental value of market and RMT features, and Set~C tests whether network topology adds further information beyond these baselines.
